\documentclass[11pt]{article}

\usepackage[T1]{fontenc}
\usepackage[utf8]{inputenc}
\usepackage[margin=1in]{geometry}
\usepackage{amsmath,amssymb,amsthm,mathtools}
\usepackage[colorlinks=true,linkcolor=blue,citecolor=blue,urlcolor=blue]{hyperref}

\numberwithin{equation}{section}
\newtheorem{theorem}{Theorem}[section]
\newtheorem{proposition}[theorem]{Proposition}
\newtheorem{lemma}[theorem]{Lemma}
\newtheorem{corollary}[theorem]{Corollary}
\newtheorem{conjecture}[theorem]{Conjecture}
\newtheorem{problem}[theorem]{Problem}
\theoremstyle{definition}
\newtheorem{definition}[theorem]{Definition}
\newtheorem{example}[theorem]{Example}
\theoremstyle{remark}
\newtheorem{remark}[theorem]{Remark}

\DeclareMathOperator{\Var}{Var}
\DeclareMathOperator{\Prob}{P}
\newcommand{\E}{\mathbb E}
\newcommand{\N}{\mathbb N}
\newcommand{\Z}{\mathbb Z}
\newcommand{\1}{\mathbf 1}

\title{Zeta-Law Entropy and Zeta-Inequality Framework}
\author{/project-pudim}
\date{}

\begin{document}
\maketitle

\begin{abstract}
We develop a probabilistic framework for organizing identities and inequalities involving the Riemann zeta function through Gibbs laws, residue distributions, and Mellin-Planck normalizations.  The basic zeta law has partition function \(\zeta(\beta)\), so zeta ratios become moments and logarithmic derivatives become energy averages and variances.  This point of view gives a self-contained derivation of Euler-score identities from the von Mangoldt divisor formula and places finite divisor sums, modular successor entropies, and Gamma-zeta log-convexity within a common notation.  The framework also isolates precise application problems concerning reciprocal-zeta curvature, a positivity question of Nantomah, and a generalized Holder inequality of Sroysang.  These applications are formulated as mathematical frontiers whose final resolution depends on completing the indicated analytic estimates, while the proved identities supply the reusable structural core.
\end{abstract}

\noindent\textbf{Keywords.} Riemann zeta function; Gibbs measure; free energy; von Mangoldt function; entropy; Mellin transform; zeta inequalities.

\section{Introduction}

Many inequalities for \(\zeta(s)\) compare neighboring zeta values, logarithmic derivatives, or Gamma-zeta products.  A useful way to normalize these expressions is to treat \(\zeta(\beta)\) as a partition function and \(n^{-\beta}\) as a Gibbs weight.  Under this normalization, zeta ratios are ordinary moments, the logarithmic derivative of \(\zeta\) is an expected energy, and convexity statements become variance inequalities.

The purpose of this paper is to make that normalization explicit and to record the verified identities that follow from it.  The main structural output is an Euler-score decomposition: the mean energy of the zeta law can be recovered either from \(-\zeta'(\beta)/\zeta(\beta)\) or from divisibility probabilities weighted by the von Mangoldt function.  We also record a finite version of the same identity and set up two related extensions: modular successor entropy and Mellin-Planck convexity.

The final section explains how the notation interfaces with three external zeta-inequality problems: a curvature question for \(1/\zeta\), a positivity problem posed by Nantomah, and a generalized Holder inequality studied by Sroysang.  Those items are stated conservatively as application problems rather than theorems.

\section{Notation and the zeta law}

\begin{definition}[Riemann zeta probability law]
For \(\beta>1\), define the probability law \(\rho_\beta\) on \(\N\) by
\begin{equation}
\label{eq:zeta-law}
\rho_\beta(n)=\frac{n^{-\beta}}{\zeta(\beta)},\qquad n\in \N.
\end{equation}
Equivalently, if \(E(n)=\log n\) and \(Z(\beta)=\zeta(\beta)\), then
\begin{equation}
\label{eq:gibbs-form}
\rho_\beta(n)=\frac{e^{-\beta E(n)}}{Z(\beta)}.
\end{equation}
\end{definition}

\begin{definition}[Zeta free energy]
The zeta free energy is
\begin{equation}
\label{eq:free-energy}
A(\beta)=\log \zeta(\beta),\qquad \beta>1.
\end{equation}
\end{definition}

\begin{proposition}[Free-energy derivatives]
For every \(\beta>1\),
\begin{equation}
\label{eq:free-derivatives}
A'(\beta)=-\E_\beta[\log N],
\qquad
A''(\beta)=\Var_\beta(\log N).
\end{equation}
\end{proposition}

\begin{proof}
Absolute convergence of the Dirichlet series and its differentiated series on compact subintervals of \((1,\infty)\) permits termwise differentiation.  Thus
\begin{equation}
\label{eq:first-derivative-proof}
A'(\beta)=\frac{\zeta'(\beta)}{\zeta(\beta)}
=-\frac{\sum_{n\ge1}(\log n)n^{-\beta}}{\zeta(\beta)}
=-\E_\beta[\log N].
\end{equation}
Differentiating once more gives the usual Gibbs identity
\begin{equation}
\label{eq:second-derivative-proof}
A''(\beta)=\E_\beta[(\log N)^2]-\E_\beta[\log N]^2
=\Var_\beta(\log N).
\end{equation}
\end{proof}

\begin{corollary}[Zeta ratios as moments]
For \(r\ge 0\) and \(\beta>1\),
\begin{equation}
\label{eq:zeta-ratio-moment}
\E_\beta[N^{-r}]=\frac{\zeta(\beta+r)}{\zeta(\beta)}.
\end{equation}
\end{corollary}

\begin{proof}
Substitute \eqref{eq:zeta-law} and sum the resulting absolutely convergent Dirichlet series:
\begin{equation}
\label{eq:zeta-ratio-proof}
\E_\beta[N^{-r}]
=\sum_{n\ge1}n^{-r}\frac{n^{-\beta}}{\zeta(\beta)}
=\frac{\zeta(\beta+r)}{\zeta(\beta)}.
\end{equation}
\end{proof}

\section{Modular successor entropy}

\begin{definition}[Residue distribution]
For \(q\ge2\), define the residue distribution of the zeta law modulo \(q\) by
\begin{equation}
\label{eq:residue-distribution}
\mu_{q,\beta}(a)=
\sum_{\substack{n\ge1\\ n\equiv a\pmod q}}\rho_\beta(n),
\qquad a\in \Z/q\Z.
\end{equation}
\end{definition}

\begin{definition}[Successor entropy]
Writing \(a+1\) cyclically in \(\Z/q\Z\), define
\begin{equation}
\label{eq:successor-entropy}
B_q(\beta)=
\sum_{a\in \Z/q\Z}
\mu_{q,\beta}(a)
\log\frac{\mu_{q,\beta}(a)}{\mu_{q,\beta}(a+1)}.
\end{equation}
\end{definition}

\begin{problem}[Successor entropy resolution]
Establish the identity
\begin{equation}
\label{eq:sigma-series}
\Sigma(\beta)
=\sum_{n=1}^{\infty}\rho_\beta(n)
\log\frac{\rho_\beta(n)}{\rho_\beta(n+1)}
=\frac{\beta}{\zeta(\beta)}
\sum_{k=1}^{\infty}\frac{(-1)^{k+1}}{k}\zeta(\beta+k),
\end{equation}
together with the modular variational formula
\begin{equation}
\label{eq:modular-variational}
\Sigma(\beta)=\sup_{q\ge2}B_q(\beta)=\lim_{q\to\infty}B_q(\beta).
\end{equation}
The expected route expands \(\log(1+1/n)\), uses the log-sum inequality, and lets large moduli isolate finite initial intervals.
\end{problem}

\begin{problem}[Prime-modulus Dirichlet \(L\)-resolution]
For prime \(p\), let \(\chi\) range over Dirichlet characters modulo \(p\), and set
\begin{equation}
\label{eq:dirichlet-character-sum}
S_a(p,\beta)=\sum_{\chi\bmod p}\chi(a)L(\beta,\chi).
\end{equation}
Prove the residue formula
\begin{equation}
\label{eq:prime-residue}
\mu_{p,\beta}(a)=\frac{S_a(p,\beta)}{(p-1)\zeta(\beta)}
\quad (1\le a\le p-1),
\qquad
\mu_{p,\beta}(0)=p^{-\beta},
\end{equation}
and derive the corresponding prime-modulus expression for \eqref{eq:successor-entropy}.
\end{problem}

\section{Euler-score identities}

\begin{theorem}[Euler-score decomposition]
For every \(\beta>1\),
\begin{equation}
\label{eq:euler-score}
-\bigl(\log\zeta\bigr)'(\beta)
=\E_\beta[\log N]
=\sum_{d=1}^{\infty}\Lambda(d)\mu_{d,\beta}(0)
=\sum_{d=1}^{\infty}\frac{\Lambda(d)}{d^\beta}.
\end{equation}
Equivalently,
\begin{equation}
\label{eq:log-zeta-dirichlet}
-\frac{\zeta'(\beta)}{\zeta(\beta)}
=\sum_{d=1}^{\infty}\frac{\Lambda(d)}{d^\beta}.
\end{equation}
\end{theorem}

\begin{proof}
Use the von Mangoldt identity
\begin{equation}
\label{eq:mangoldt-identity}
\log n=\sum_{d\mid n}\Lambda(d).
\end{equation}
Then, by nonnegative summation and \eqref{eq:zeta-law},
\begin{equation}
\label{eq:euler-score-proof}
\E_\beta[\log N]
=\sum_{n\ge1}\rho_\beta(n)\sum_{d\mid n}\Lambda(d)
=\sum_{d\ge1}\Lambda(d)\Prob_\beta(d\mid N).
\end{equation}
The divisibility probability is
\begin{equation}
\label{eq:divisibility-probability}
\Prob_\beta(d\mid N)
=\frac{1}{\zeta(\beta)}\sum_{m\ge1}(dm)^{-\beta}
=d^{-\beta}.
\end{equation}
Substituting \eqref{eq:divisibility-probability} into \eqref{eq:euler-score-proof} gives \eqref{eq:euler-score}; \eqref{eq:log-zeta-dirichlet} follows from \eqref{eq:first-derivative-proof}.
\end{proof}

\begin{proposition}[Finite Euler-score identity]
For every \(A\subseteq\{1,\ldots,N\}\),
\begin{equation}
\label{eq:finite-euler-score}
\sum_{n\in A}\log n
=\sum_{d\le N}\Lambda(d)\sum_{m\le N/d}\1_A(md).
\end{equation}
\end{proposition}

\begin{proof}
Reverse the order of summation:
\begin{equation}
\label{eq:finite-euler-score-proof}
\sum_{d\le N}\Lambda(d)\sum_{m\le N/d}\1_A(md)
=\sum_{n\in A}\sum_{d\mid n}\Lambda(d)
=\sum_{n\in A}\log n,
\end{equation}
where the last equality is \eqref{eq:mangoldt-identity}.
\end{proof}

\section{Mellin-Planck normalization}

\begin{definition}[Mellin-Planck partition function]
For \(s>1\), define
\begin{equation}
\label{eq:mellin-planck}
M(s)=\Gamma(s)\zeta(s)
=\int_0^\infty \frac{t^{s-1}}{e^t-1}\,dt.
\end{equation}
The normalized law associated with \(M\) is
\begin{equation}
\label{eq:mellin-planck-law}
\nu_s(dt)=\frac{t^{s-1}}{(e^t-1)M(s)}\,dt.
\end{equation}
\end{definition}

\begin{proposition}[Log-convexity of the Mellin-Planck partition function]
For \(s>1\),
\begin{equation}
\label{eq:mellin-logconvex}
\frac{d^2}{ds^2}\log M(s)=\Var_{\nu_s}(\log t)\ge 0.
\end{equation}
In particular, \(M\) is log-convex on \((1,\infty)\).
\end{proposition}

\begin{proof}
Differentiate the logarithm of \eqref{eq:mellin-planck} under the integral sign.  The first derivative is the \(\nu_s\)-expectation of \(\log t\), and the second derivative is the corresponding variance, giving \eqref{eq:mellin-logconvex}.
\end{proof}

\section{Application problems}

\begin{problem}[Nantomah positivity]
Determine whether, for every \(n\in\N\),
\begin{equation}
\label{eq:nantomah-positivity}
(n+2)\zeta(n+1)\zeta(n+3)
-(n+1)\zeta(n+2)^2
-\zeta(n+1)\zeta(n+2)>0.
\end{equation}
The zeta-law approach rewrites the expression using moments of \(X(m)=1/m\) under \(\rho_s\), then seeks a positive linear part and a controlled quadratic correction.
\end{problem}

\begin{problem}[Alzer-Kwong reciprocal-zeta curvature]
Let \(F(x)=1/\zeta(x)\).  Prove or refute the sign pattern
\begin{equation}
\label{eq:alzer-kwong-curvature}
F''(x)>0\quad\text{on }(-4n,-4n+2),
\qquad
F''(x)<0\quad\text{on }(-4n-2,-4n)
\end{equation}
for every integer \(n\ge1\).  A natural route uses the zeta functional equation and positive-axis curvature bounds for \((\log\zeta)''(s)\).
\end{problem}

\begin{problem}[Sroysang generalized Holder inequality]
Let \(m\ge2\), \(r\ge1\), \(x_1,\ldots,x_m>1\), and \(p_1,\ldots,p_m>1\) satisfy
\begin{equation}
\label{eq:sroysang-condition}
\sum_{i=1}^m\frac1{p_i}=\frac1r.
\end{equation}
With
\begin{equation}
\label{eq:sroysang-t}
T=r\sum_{i=1}^m\frac{x_i}{p_i},
\end{equation}
establish conditions under which
\begin{equation}
\label{eq:sroysang-inequality}
\zeta(T)\le
\frac{\prod_{i=1}^m(\Gamma(x_i)\zeta(x_i))^{r/p_i}}{\Gamma(T)}.
\end{equation}
The Mellin-Planck formulation suggests applying generalized Holder to
\begin{equation}
\label{eq:sroysang-holder-functions}
f_i(t)=\left(\frac{t^{x_i-1}}{e^t-1}\right)^{1/p_i}.
\end{equation}
\end{problem}

\section{Conclusion}

The zeta-law normalization turns several analytic expressions into standard probabilistic objects.  Equations \eqref{eq:zeta-ratio-moment}, \eqref{eq:free-derivatives}, \eqref{eq:euler-score}, and \eqref{eq:mellin-logconvex} are the proved structural core: zeta ratios are moments, logarithmic derivatives are energy averages, Euler scores are divisibility averages, and Mellin-Planck convexity is a variance identity.  The modular entropy formulas \eqref{eq:sigma-series} and \eqref{eq:modular-variational}, together with the application problems \eqref{eq:nantomah-positivity}, \eqref{eq:alzer-kwong-curvature}, and \eqref{eq:sroysang-inequality}, define the next analytic frontier.

\begin{thebibliography}{9}

\bibitem{alzer-kwong}
H. Alzer and M. K. Kwong,
\emph{On the concavity and convexity of \(1/\zeta\)},
International Journal of Number Theory 21, no. 8 (2025), 1825--1835.
\url{https://doi.org/10.1142/S1793042125500897}.

\bibitem{nantomah}
K. Nantomah,
\emph{Open Problem on Riemann Zeta Function},
ResearchGate problem note, October 2024.
\url{https://www.researchgate.net/publication/384676538_Open_Problem_on_Riemann_Zeta_Function}.

\bibitem{sroysang}
B. Sroysang,
\emph{Two Inequalities for the Riemann Zeta Functions},
Mathematica Aeterna 3, no. 1 (2013), 21--24.
\url{https://arastirmax.com/en/system/files/dergiler/135290/makaleler/3/1/arastirmax-two-inequalities-riemann-zeta-functions.pdf}.

\end{thebibliography}

\end{document}
